\begin{statement}{P3F.38}{Proved}
For the OS50 source \(I_1^*+IV+I_1\), use
\(M=\mathcal O(aG-O-L_0)\) and the simple whole zero at its
\(I_1\) value. The following actual two-cap constructions exist:
\[
\begin{array}{c|c|c}
(a,b)&\text{star cap}&\text{IV cap}\\\hline
(6,2)&(2,n_*),\ n_*\text{ odd}&(m,n_4),\ m\equiv1\pmod3\\
(2,6)&(2,n_*),\ n_*\text{ odd}&(m,n_4),\ 3\mid m
\end{array}
\]
Here \(m>0\), \(\gcd(m,n_4)=1\), and the actual elliptic,
gamma and delta collars are retained. Put
\(D=mn_*+2n_4\). Their circle orbits have
\[
 \pi_1(Y)=\mathbb Z/|D|,\qquad
 H_j(Y;\mathbb Z)=
 \begin{cases}\mathbb Z,&j=0,5,\\
 \mathbb Z/|D|,&j=1,3,\\0,&j=2,4.
 \end{cases}
\]
The actual six-manifold is a smooth sphere exactly when
\(D=1\) or \(D=-1\). Equivalently, in the first row
\(m\equiv1\pmod6\), and in the second row \(m=3k\) with
\(k\) odd; in both rows
\[
 n_*\text{ is odd},\qquad
             n_4=(\varepsilon-mn_*)/2,
                    \quad\varepsilon\in\{1,-1\}.
\]
The case \(m=1,n_4=0\) in the first row agrees with P3F.36.
No parameter-return class is assigned by this statement.
\end{statement}

\begin{proof}
\par\medskip\noindent\textbf{The two actual IV roots.}
We use the marked reduced-IV equation
\[
 y^2=x^3+q^2a(q)x+q^2b(q),\qquad b(0)\ne0,
\]
where these coefficient functions are distinct from the source
integers \((a,b)\). For \((a,b)=(6,2)\), put
\(m=3\ell+1\), \(q=z^m\), \(x=z^{2\ell}X\),
\(y=z^{3\ell}Y\). The exact equation becomes
\[
 Y^2=X^3+z^{2\ell+2}a(z^m)X+z^2b(z^m).
\]
Its minimal regular central fibre is IV. The deck is
\[
 (X,Y,z)\longmapsto
 (\zeta_m^{-2\ell}X,\zeta_m^{-3\ell}Y,\zeta_mz).
\]
In the blowup chart \(X=zX_1,Y=zY_1\), the two nonidentity
components are labelled by \(Y_1=\pm\sqrt{b(0)}\).
The character on \(Y_1\) is \(\zeta_m^{-(3\ell+1)}=1\),
so the deck fixes each component label. The root line
\(N=\mathcal O(6G'-O')\) has degree zero on each component,
and \(Q'=2G'\) has component two. Multiplication by \(m\)
is an analytic local isomorphism of the smooth Neron group:
it has derivative \(m\) on the additive identity component and
is invertible on the component group \(\mathbb Z/3\).
The selected local section \(R\) with \(mR=Q'\) is therefore
unique and deck-invariant.

The root line is compared with the actual source line on the
resolved fibre product. Every curve contracted in relative
minimalization has degree zero for the pulled-back source line,
since its source component degree is zero. Descending through the
actual point contractions gives the root line. Its ratio with
\(\mathcal O(6G'-O')\) is generically trivial and has degree
zero on every component, so the proper-direct-image construction
of P3K.15 gives an integral whole-fibre divisor for that ratio.
Multiplication by the corresponding base power supplies the
actual line isomorphism. This retains its deck character.
The source correction \(-L_0\) is supported at the other,
starred value and has no local term here.

For \((a,b)=(2,6)\), put \(m=3\ell\),
\(q=z^m\), \(x=z^{2\ell}X\), \(y=z^{3\ell}Y\).
The exact root equation is
\[
 Y^2=X^3+z^{2\ell}a(z^m)X+b(z^m).
\]
Its central curve is the smooth elliptic curve
\(E_0:Y^2=X^3+b(0)\). The deck on it is
\((X,Y)\mapsto(\zeta_3X,Y)\). Its fixed points are
\(P_0=O,P_1,P_2\), where \(P_1,P_2\) are the two nonzero
three-torsion points. The original marked nonidentity component
of \(G\) identifies \(G'_0=P_1\), and \(2G'_0=P_2\).
Thus \(Q'=6G'\) has central value zero. The unique analytic
\(m\)-division \(R\) of \(Q'\) with \(R_0=0\) is
deck-invariant.

\par\medskip\noindent\textbf{The good-root line and its exact exceptional correction.}
First take the minimal cubic root \(q=s^3\). At each of the
three fixed points of
\((X,Y,s)\mapsto(\zeta X,Y,\zeta s)\), the two tangent
weights are \((1,1)\). Blow up these points by
\(\beta:\widehat S\to S_{\rm root}\), with exceptional
curves \(E_i\). In coordinates \((w/s,s)\), the quotient
coordinates are \((w/s,s^3)\). Hence the quotient
\(\pi:\widehat S\to\widetilde S\) is smooth and
\(\pi^*L_i=3E_i\). Since \(E_i^2=-1\), its image has
\(L_i^2=-3\).

The proper central elliptic curve has square \(-3\) after the
three blowups. Its quotient image \(C\) has square \(-1\),
because its pullback is that proper transform and the quotient
has degree three. It meets the three \(L_i\) at distinct points.
The actual smooth point contraction
\(\alpha:\widetilde S\to S\) contracts \(C\). The three
images of \(L_i\) then have square \(-2\) and meet at one
point with distinct tangent directions, giving the minimal regular
IV fibre. The invariant functions
\(q=s^3,x=s^2X,y=s^3Y\) identify the displayed Weierstrass
model off the exceptional curves. The resolved graph and relative
minimality identify this model with the original source over the
disc.

The section curves avoid \(C\): on the root they are transverse
to the base and their proper transforms meet \(E_i\) away from
the proper central curve. Pullback of the actual two section
divisors therefore gives
\[
 \boxed{\ \beta^*\mathcal O(2G'-O')
       =\pi^*\alpha^*\mathcal O(2G-O)+E_2-E_0\ }.             \tag{P3F.38.1}
\]
The deck weights of this line at \((P_0,P_1,P_2)\) are
\((\zeta^{-1},1,\zeta)\), up to one common character. Indeed
for \(\mathcal O(P_2)\) the canonical section is a local
parameter of weight \(\zeta\) times its frame, so the frame
has weight \(\zeta\); the inverse divisor at \(P_0\) gives
\(\zeta^{-1}\). These are exactly the weights of the
correction in (P3F.38.1).

Translation by \(Q'=6G'\) fixes all three central points and
preserves \(E_2-E_0\). Its canonical divisor lift combined with
the pulled-back source lift therefore acts on the right side of
(P3F.38.1), and descends through \(\beta\) to the actual map
on \(N\). This agrees with the original punctured lift and
commutes with the deck. For \(m=3\ell\), pull this entire
correspondence through \(s=z^\ell\) and resolve the fibre
product. The line identity pulls back exactly; all extra
exceptional terms are its actual pullbacks.

The nonzero fixed central Picard point gives the generator of the
order-three Euler torsion of the cap. Explicitly, for the Picard
torus \(V/\Lambda\), the connecting map is
\[
 (V/\Lambda)^D\longrightarrow\Lambda/(D-1)\Lambda,
                  \qquad[v]\longmapsto[(D-1)v].
\]
It is an isomorphism because \(D-1\) is invertible on \(V\),
and its target has order three. The line point \(P_2\) is
nonzero, so its clutch is the nonzero torsion class. A common
deck character has zero absolute first Chern class but remains
in the actual line lift.

\par\medskip\noindent\textbf{The commuting norm lifts and free quotients.}
For either IV root, first lift \(t_{-R}\) on \(N\), before
taking its \(n_4\)-th power. The actual meromorphic chord section
of its Hom line has an integral vertical correction. The
component-degree equations make that correction a whole fibre;
multiplying by its base power gives an invertible holomorphic
line map \(F\). Average in the deck-conjugation action on the
rank-one Hom module, choosing its nonzero central character.
The resulting lift is still invertible after shrinking and satisfies
\[
                   DFD^{-1}=\chi F,\qquad\chi^m=1.
\]
The maps \(F^m\) and the specified \(B_0\) cover the same
translation \(t_{-Q'}\), so
\(B_0=\beta(z)F^m\) for a base unit. Its actual descent gives
\(\beta(\zeta_mz)=\beta(z)\). Thus \(F,D\) both commute
with \(B_0\), and
\[
 (F^{n_4}D)^m/B_0^{n_4}
              =\chi^{n_4m(m-1)/2}\beta(q)^{-n_4}.              \tag{P3F.38.2}
\]
This is an invariant nonvanishing unit. Divide by its analytic
\(m\)-th root to obtain \(A_0\), retaining the displayed
character factor, so
\(A_0^m=B_0^{n_4}\) and \([A_0,B_0]=1\).
The source map is \(t_{-n_4R}D\).

Combine this with the actual star construction of P3F.36. For
\(L=\operatorname{lcm}(2,m)\), use the parameter map
\(\lambda=r^L\) and multiply the star and IV lifts respectively
by \(r^{Ln_*/2}\) and \(r^{Ln_4/m}\). Their exact norm
identities hold for \(B_\lambda\). The primitive generators
from Bezout have positive scalar powers \(r^{L/2}\) and
\(r^{L/m}\); their remaining maps and inverses have bounded
norm ratios on compact smaller root tubes. They are therefore
uniformly contracting for small \(r\ne0\). Their quotients
are smooth, proper and compact over the original discs by the
actual norm shells. A unit line rotation cannot equal a nonzero
iterate of such a contraction, so the delta action is free there.
The line correspondences identify their punctured collars with
the actual exterior. This constructs the stated global families
under the displayed parameter map, rather than inferring them
from their homology.

The positive angular convention is the same at both caps. On a
root of order \(m_i\), write \(A_i:(\theta,t)\mapsto
(\theta+1/m_i,t-n_i/m_i)\), with \(B:t\mapsto t-1\).
The original coordinates are \(x=m_i\theta,y=t+n_i\theta\).
Thus \(\alpha_i\mapsto-n_ih_i\),
\(\gamma\mapsto m_ih_i\), and the killed meridian is exactly
\(m_i\alpha_i+n_i\gamma\).

\par\medskip\noindent\textbf{The actual common source complement.}
Remove both additive source tubes and call the remaining surface
\(P\). Its fundamental group is exactly the base annulus group.
To see this without losing an elliptic kernel, let the starred
monodromy be \(-T_u\), and let \(T_v\) be the monodromy of
the remaining \(I_1\), for primitive cycles \(u,v\). The IV
monodromy has trace \(-1\). The product relation and transvection
formula give
\[
 -1=-\operatorname{tr}(T_u^{-1}T_v^{-1})
                         =-2+(u\cdot v)^2.
\]
Hence \(|u\cdot v|=1\). The nodal vanishing disc kills \(v\),
and conjugation by the retained starred meridian kills
\((-T_u)v=-v\pm u\). Thus the elliptic lattice dies, leaving
\(\pi_1(P)=\mathbb Z\alpha\), with
\(\alpha_4=\alpha_*^{-1}\). Translation fixes this generator,
so \(\pi_1(W)=\mathbb Z\alpha\oplus\mathbb Z\gamma\).

We compute the homology from actual component inclusions. In the
NS basis of P3F.36, the form \(B\) is fixed by
\(O^2=G^2=-1\), \(O\cdot F=G\cdot F=1\), the negative
\(D_5+A_2\) root matrix, and the intersections of \(G\) with
\(R_0,A_{2,1}\); the other pairings are zero. Its determinant
is \(-1\). Put
\[
 L_*=(L_0,L_1,R_0,R_1,C_0,C_1),\quad
 L_4=(F-A_{2,1}-A_{2,2},A_{2,1},A_{2,2}),\quad L=(L_*,L_4),
\]
where \(L_0=F-2C_0-2C_1-L_1-R_0-R_1\).
Both local fibre neighbourhoods have free homology only in degrees
zero and two. The cone of \(C(S)\to C(S,P)\), shifted once,
therefore gives the actual complement complex
\[
 C_4=\mathbb Z,\ C_3=\mathbb Z^2,\ C_2=\mathbb Z^{10},
 \ C_1=\mathbb Z^9,\ C_0=\mathbb Z,
 \qquad d_4=(1,1)^t,\quad d_2=L^tB.                         \tag{P3F.38.3}
\]
All other differentials are zero. The two complexes before taking
this cone are even and free, so their degree-one homotopy group
is zero. This makes the construction functorial for the actual
translation, with no omitted annular term.

For \(b=6\), the actual section intersections are
\(O\cdot6G=G\cdot6G=1\), \(O\cdot7G=2\),
\(G\cdot7G=1\). The sections meet respectively \(L_1\)
with the identity IV component, and \(R_1\) with \(A_{2,1}\).
Solving against \(B\) determines their integral vectors and
hence the translation matrix. For \(b=2\) use the actual
\(2G,3G\) classes of P3F.36. In each case invert the positive
translation matrix to obtain the contracting map \(t_{-bG}\).
The certificates check the isometry of \(B\), all component
permutations, and the equivariance of (P3F.38.3).

For each boundary use the complex with sole differential
\(L_i^tBL_i:\mathbb Z^{k_i}\to\mathbb Z^{k_i}\).
Its map into (P3F.38.3) is \(L_i\) in degree two, the coordinate
inclusion into \(\mathbb Z^9\) in degree one, and the appropriate
coordinate of \(\mathbb Z^2\) in degree three. Taking the actual
mapping-torus cones gives every boundary map. We obtain
\[
 H_1(W)=\mathbb Z^2,\quad H_3(W)=\mathbb Z^2,\quad H_4(W)=\mathbb Z,
 \qquad
 H_2(W)=\begin{cases}\mathbb Z^2\oplus\mathbb Z/2,&b=2,\\
                         \mathbb Z^2\oplus\mathbb Z/6,&b=6.
          \end{cases}
\]

\par\medskip\noindent\textbf{Every integral degree-two map.}
For \(b=6\), the boundary bases are
\((\epsilon_*,r_*,s_*)\), of orders \((2,0,0)\), and
\((\epsilon_4,r_4,s_4)\), of orders \((3,0,0)\).
In the exterior basis \((\tau,g_1,g_2)\), of orders
\((6,0,0)\), the actual matrix is
\[
 \begin{pmatrix}
 -15&-5&83&-4&5&-4\\0&0&-1&0&0&0\\0&0&-2&0&0&1
 \end{pmatrix}.
\]
The fibre classes are \(F_*=-\epsilon_*+2r_*\) and
\(F_4=r_4\), both mapping to \(5\tau\). The star fibre
maps to \(2C_*\), and the good-root fibre maps to \(C_4\)
with degree one. Thus the cap rows are
\((0,1,\kappa_*)\) and \((0,1,\kappa_4)\), for their actual
integral collar coefficients. The torsion map sends
\(\epsilon_*\) to \(3\tau\) and \(\epsilon_4\) to
\(2\tau\). The exact Chinese-remainder map
\(\mathbb Z/6\to\mathbb Z/2\oplus\mathbb Z/3\) is
\(1\mapsto(1,-1)\). In this single torsion generator followed
by the four free generators, the complete map is
\[
 D_2^{\rm good}=\begin{pmatrix}
 1&1&5&5&2\\0&0&-1&0&0\\0&0&-2&0&1\\
 0&-1&-\kappa_*&0&0\\0&0&0&-1&-\kappa_4
 \end{pmatrix}.                                               \tag{P3F.38.4}
\]
Its first row is read modulo six; these are the exact preceding
entries reduced by \(6\tau=0\). Its determinant is one.

For \(b=2\), the star boundary basis is the same, while the IV
boundary has two free generators \((r_4,s_4)\). The exterior
basis has orders \((2,0,0)\), and its matrix is
\[
 \begin{pmatrix}-5&2&33&3&9\\0&0&1&0&-2\\0&0&-2&0&3\end{pmatrix}.
\]
Here \(F_*=-\epsilon_*+2r_*\), \(F_4=3r_4\). The reduced
IV root has three rational components permuted transitively, so
its fibre maps to \(3C_4\). Its cap row is again
\((1,\kappa_4)\). Thus the complete matrix, with its first row
read modulo two, is
\[
 D_2^{\rm red}=\begin{pmatrix}
 1&0&1&1&1\\0&0&1&0&-2\\0&0&-2&0&3\\
 0&-1&-\kappa_*&0&0\\0&0&0&-1&-\kappa_4
 \end{pmatrix}.                                               \tag{P3F.38.5}
\]
Its free core block has determinant
\(\det\left(\begin{smallmatrix}1&-2\\-2&3\end{smallmatrix}\right)=-1\),
and the full map is unimodular. Both maps are therefore integral
isomorphisms for their actual collars. No source inversion has
been extended across a marked cap.

\par\medskip\noindent\textbf{The other maps and the specialization indices.}
The good-root generic elliptic fibre transports through its smooth
root family to \(E_0\) with degree one. On \(H_1\), that
same transport is the identity lattice \(\Lambda\). If
\(u n_4+vm=1\), the cap monodromy is \(D^u\). The two
boundary and cap torsion sublattices are equal integrally:
\[
 (D^u-1)\Lambda=(D-1)\Lambda,
\]
because \(u\) is a unit modulo three and, for \(u\equiv2\),
\(D^2-1=(D-1)(D+1)\) with \(D+1=-D^2\) unimodular.
Thus the actual boundary \(\mathbb Z/3\) maps isomorphically
to cap torsion. The starred \(\mathbb Z/2\) map is the
primitive graph specialization of P3F.36. For the reduced IV cap
both groups have no degree-one torsion.

Retain all path coefficients in the free columns. The preceding
torsion isomorphisms eliminate them without changing any free
relation. In the positive meridian basis
\((\alpha_*,\gamma_*,\alpha_4,\gamma_4)\), the free
Mayer--Vietoris matrix is
\[
 D_1=\begin{pmatrix}
 1&0&-1&0\\0&1&0&1\\n_*&-2&0&0\\0&0&n_4&-m
 \end{pmatrix}.                                               \tag{P3F.38.6}
\]
Its determinant is \(-D\). In degree three use the oriented
source boundaries and their positive fibre--gamma traces. The
exterior rows are \((0,1,0,1)\) and \((-1,0,1,0)\); in
the reduced case reverse the certificate's second target generator
to obtain this same convention. The cap rows are
\((-n_*,2)\) and \((-n_4,m)\), by the actual norm powers.
Graph corrections lie inside the two-dimensional central source
and add no three-class. Hence
\[
 D_3=\begin{pmatrix}
 0&1&0&1\\-1&0&1&0\\n_*&-2&0&0\\0&0&n_4&-m
 \end{pmatrix},\qquad\det D_3=-D.                           \tag{P3F.38.7}
\]
The degree-four map is \((1,-1)\), and the degree-zero map is
\(\left(\begin{smallmatrix}1&1\\-1&0\\0&-1\end{smallmatrix}\right)\).
These specify every nonzero map. The chain matrices, representatives,
source meridians and oriented fibre traces are executed in
\path{code/part3/i1star_iv_boundary.py} and
\path{code/part3/i1star_iv_reduced_boundary.py}, with their JSON
outputs.

For van Kampen, all elliptic periods die in \(W\). Their cap
images generate the star graph group and, in the good case, the
whole elliptic lattice. The reduced-IV central curve is simply
connected. After killing these normal subgroups, the exact remaining
presentation is
\[
 \langle h_*,h_4\mid h_*^2=h_4^m,
                         \ h_*^{n_*}=h_4^{-n_4}\rangle.
\]
Choose \(r_0n_*+2s_0=1\). Both right-hand powers are powers of
\(h_4\), so \(h_*=h_4^{-r_0n_4+s_0m}\). Thus the group is
cyclic, and the two exact relations give \(\mathbb Z/|D|\).
The matrices give the asserted integral homology. Here \(D\ne0\):
otherwise primitivity would imply \(m\mid2\), while the allowed
\(m\) and odd \(n_*\) exclude zero determinant also at \(m=1\).

For an arbitrary allowed one-vertex cusp fan, use the finite sequence
of actual periodic diagonal flops of P3F.35. Each step is the common
blowup of its embedded normal \(\mathcal O(-1)\oplus\mathcal O(-1)\)
curve followed by the other ruling contraction. These maps are the
identity outside the cusp, so they preserve the present marked caps.
Their integral blowup decompositions and van Kampen maps preserve
homology and fundamental group. Thus the sphere question is exactly
the same as for the vertical fan used in the orbit calculation.

Finally, the unchanged whole-zero \(I_1\) has the same actual
fixed sphere and primitive Hopf meridian as P3F.36. Its Gysin and
fixed-tube gluing gives the smooth six-sphere when \(|D|=1\),
and gives \(\pi_1(X)=\pi_1(Y)\) for every tuple, excluding
all other determinants. The stated arithmetic form follows because
\(|D|=1\) forces odd \(m\), and then
\(2n_4\equiv\varepsilon\pmod m\) automatically gives
\(\gcd(m,n_4)=1\). The actual line and norm lifts, rather than
only this arithmetic, were constructed above.
\end{proof}
